% ============================================================================
%  ESTADÍSTICA INFERENCIAL — Sesión 10
%  Semana 10: 12 al 17 de octubre de 2026   (lunes 12: festivo)
%  2.6 Estimación puntual · 2.7 Conceptos y tipos de intervalos de confianza
% ============================================================================
\documentclass[aspectratio=169,10pt]{beamer}

\usepackage[utf8]{inputenc}
\IfFileExists{spanish.ldf}{\usepackage[spanish,es-noshorthands]{babel}}{}
\usepackage{amsmath,amssymb}
\usepackage{booktabs}
\usepackage{tikz}
\usepackage{listings}

\usetheme{Uniminuto}

\uniasignatura{Estadística inferencial}
\uniprofesor{Andrés Felipe Puerta Velez}
\unicorreo{andres.puerta-v@uniminuto.edu.co}
\uninrc{CBASBL061}
\unisemana{Sesión 10 — Semana 10 · 12 al 17 de octubre de 2026}

\newcommand{\listatiny}{\setbeamerfont{itemize/enumerate body}{size=\tiny}\tiny}

\newcommand{\fx}[1]{{\ttfamily\color{uniAzul}#1}}

% ------------------------------------------------------------- CÓDIGO R ----
\lstdefinestyle{Rcode}{
  language=R,
  basicstyle=\ttfamily\scriptsize,
  keywordstyle=\color{uniAzul}\bfseries,
  commentstyle=\color{uniGris}\itshape,
  stringstyle=\color{uniTitulo},
  showstringspaces=false, numbers=none, breaklines=true,
  columns=fullflexible,
  backgroundcolor=\color{uniCremaClaro},
  frame=leftline, framerule=1.2pt, rulecolor=\color{uniNaranja},
  xleftmargin=0.35cm, framexleftmargin=0.35cm,
  framexrightmargin=0.15cm, framextopmargin=2pt, framexbottommargin=2pt,
  aboveskip=0.15cm, belowskip=0.15cm
}
\lstdefinestyle{Rmini}{
  language=R,
  basicstyle=\ttfamily\tiny,
  keywordstyle=\color{uniAzul}\bfseries,
  commentstyle=\color{uniGris}\itshape,
  stringstyle=\color{uniTitulo},
  showstringspaces=false, numbers=none, breaklines=true,
  columns=fullflexible,
  backgroundcolor=\color{uniCremaClaro},
  frame=leftline, framerule=1.2pt, rulecolor=\color{uniNaranja},
  xleftmargin=0.35cm, framexleftmargin=0.35cm,
  framexrightmargin=0.15cm, framextopmargin=2pt, framexbottommargin=2pt,
  aboveskip=0.12cm, belowskip=0.12cm
}
\lstdefinestyle{Rout}{
  basicstyle=\ttfamily\tiny\color{black},
  showstringspaces=false, numbers=none, breaklines=true,
  columns=fullflexible,
  backgroundcolor=\color{uniGris!10},
  frame=leftline, framerule=1.2pt, rulecolor=\color{uniGris},
  xleftmargin=0.35cm, framexleftmargin=0.35cm,
  framexrightmargin=0.15cm, framextopmargin=2pt, framexbottommargin=2pt,
  aboveskip=0.12cm, belowskip=0.12cm
}

\begin{document}

\portada

\begin{frame}{Orden del día}
  \begin{enumerate}
    \item Llamado a lista y revisión de los ejercicios propuestos
    \item Repaso: dónde quedamos
    \item \azul{Tema 2.6: estimación puntual}
    \item \azul{Tema 2.7: conceptos y tipos de intervalos de confianza}
    \item Ejercicios en clase
    \item Cierre y ejercicios propuestos
  \end{enumerate}
\end{frame}

\begin{frame}{Repaso: dónde quedamos}
  En la sesión 9 aprendimos que $\bar{X}$ tiene su propia distribución:
  \[ \bar{X}\sim N\!\left(\mu,\ \frac{\sigma}{\sqrt{n}}\right)
     \quad\text{(por el TLC)} \]

  \vspace{0.2cm}
  \begin{alertblock}{El giro de hoy}
    \footnotesize
    Hasta ahora conocíamos $\mu$ y preguntábamos por $\bar{x}$. \destacar{Desde
    hoy invertimos la pregunta:} observamos $\bar{x}$ en una muestra y queremos
    decir algo sobre el $\mu$ que no conocemos.

    \vspace{0.15cm}
    Eso es \azul{estimar}, y es el corazón de la inferencia estadística.
  \end{alertblock}
\end{frame}

\seccion{2.6 Estimación puntual}

\begin{frame}{Estimación puntual}
  \begin{block}{Definición}
    \small
    Un \azul{estimador puntual} es un estadístico que se usa como aproximación
    de un parámetro desconocido. Da \destacar{un solo número}.
  \end{block}

  \vspace{0.2cm}
  \footnotesize
  \begin{center}
  \begin{tabular}{@{}lll@{}}
    \toprule
    \textbf{Parámetro (población)} & \textbf{Estimador (muestra)} &
    \textbf{Se calcula} \\
    \midrule
    Media $\mu$        & $\bar{x}$   & $\sum x_i/n$ \\
    Proporción $p$     & $\hat{p}$   & $x/n$ \\
    Varianza $\sigma^{2}$ & $s^{2}$  & $\sum(x_i-\bar x)^{2}/(n-1)$ \\
    Desviación $\sigma$ & $s$        & $\sqrt{s^{2}}$ \\
    \bottomrule
  \end{tabular}
  \end{center}

  \vspace{0.2cm}
  {\footnotesize\color{uniGris} El $n-1$ de la varianza muestral no es un
  capricho: hace que $s^{2}$ sea insesgado. Con $n$ en el denominador se
  subestimaría sistemáticamente la varianza.}
\end{frame}

\begin{frame}{Propiedades de un buen estimador}
  \begin{columns}[T,onlytextwidth]
    \begin{column}{0.48\textwidth}
      \begin{block}{Insesgado}
        \footnotesize
        Su valor esperado es el parámetro: $E(\bar X)=\mu$. En promedio da en
        el blanco, sin desviarse hacia arriba ni hacia abajo.
      \end{block}
      \begin{block}{Eficiente}
        \footnotesize
        Entre varios insesgados, el de \azul{menor varianza}. Da resultados más
        parecidos entre muestras.
      \end{block}
    \end{column}
    \begin{column}{0.48\textwidth}
      \begin{block}{Consistente}
        \footnotesize
        Se acerca al parámetro a medida que $n$ crece: $\sigma/\sqrt{n}\to 0$.
      \end{block}
      \begin{alertblock}{La limitación de la estimación puntual}
        \footnotesize
        $\bar{x}=52$ es casi con seguridad \destacar{distinto} de $\mu$. Un
        solo número no dice qué tan cerca puede estar. \azul{Por eso se
        necesitan los intervalos.}
      \end{alertblock}
    \end{column}
  \end{columns}
\end{frame}

\seccion{2.7 Intervalos de confianza}

\begin{frame}{Estimación por intervalo}
  \begin{block}{Idea y estructura}
    \small
    En lugar de un número, se reporta un \destacar{rango} de valores creíbles
    para el parámetro, con un \destacar{nivel de confianza}:
    \[ \text{estimador} \;\pm\; \text{margen de error}
       \qquad\text{es decir}\qquad
       \bar{x} \;\pm\; z\cdot\frac{\sigma}{\sqrt{n}} \]
  \end{block}

  \vspace{0.2cm}
  \begin{alertblock}{Las tres piezas del margen}
    \footnotesize
    El \azul{valor crítico} ($z$ o $t$) lo fija la confianza; la
    \azul{variabilidad} ($\sigma$ o $s$) viene del fenómeno; el \azul{tamaño de
    muestra} ($n$) lo decide el investigador. Solo la tercera está bajo nuestro
    control.
  \end{alertblock}
\end{frame}

\begin{frame}{Qué significa ``95\% de confianza''}
  \begin{alertblock}{Interpretación correcta}
    \footnotesize
    Si se repitiera el muestreo muchas veces y en cada muestra se construyera
    su intervalo, cerca del \destacar{95\% de esos intervalos contendría el
    verdadero $\mu$}. La confianza está en el \azul{método}, no en el intervalo
    particular que obtuvimos.
  \end{alertblock}

  \vspace{0.2cm}
  \begin{block}{Interpretación incorrecta, y muy común}
    \footnotesize
    ``Hay 95\% de probabilidad de que $\mu$ esté entre 50 y 54''. No: $\mu$
    es un número fijo, no una variable aleatoria. O está en el intervalo o no
    está. Lo aleatorio es el intervalo, que cambia con cada muestra.
  \end{block}

  \vspace{0.15cm}
  {\footnotesize\color{uniGris} En el examen esta distinción se pregunta.}
\end{frame}

\begin{frame}{Tipos de intervalo de confianza}
  \footnotesize
  \begin{center}
  \begin{tabular}{@{}lll@{}}
    \toprule
    \textbf{Parámetro} & \textbf{Condición} & \textbf{Intervalo} \\
    \midrule
    Media $\mu$ & $\sigma$ conocida &
      $\bar{x}\pm z_{\alpha/2}\dfrac{\sigma}{\sqrt{n}}$ \\[0.5em]
    Media $\mu$ & $\sigma$ desconocida &
      $\bar{x}\pm t_{\alpha/2,\,n-1}\dfrac{s}{\sqrt{n}}$ \\[0.5em]
    Proporción $p$ & $n\hat p\ge 5$, $n(1-\hat p)\ge 5$ &
      $\hat{p}\pm z_{\alpha/2}\sqrt{\dfrac{\hat p(1-\hat p)}{n}}$ \\
    \bottomrule
  \end{tabular}
  \end{center}

  \vspace{0.2cm}
  \begin{exampleblock}{Valores críticos más usados}
    \footnotesize
    \centering
    90\%: $z=1{,}645$ \qquad 95\%: $z=1{,}96$ \qquad 99\%: $z=2{,}576$
  \end{exampleblock}
\end{frame}

\begin{frame}{Ejercicio 1 — en clase}
  \begin{alertblock}{Individual (9 minutos)}
    Una muestra de \destacar{$n=64$} facturas da un valor promedio de
    \destacar{$\bar{x}=52$} (miles de pesos). Por registros históricos se sabe
    que \destacar{$\sigma=8$}.

    \vspace{0.25cm}
    \begin{enumerate}
      \item Construya el intervalo de confianza del 95\% para $\mu$.
      \item Constrúyalo al 99\%.
      \item ¿Cuál es más ancho y por qué?
      \item Interprete el intervalo del 95\% en una frase.
    \end{enumerate}
  \end{alertblock}
\end{frame}

\begin{frame}{Ejercicio 1 — solución}
  \small
  \begin{block}{Cálculos}
    \[ \sigma_{\bar x}=\frac{8}{\sqrt{64}}=1{,}00 \]
    \begin{align*}
      95\%:\;& 52\pm 1{,}96(1{,}00)=52\pm 1{,}96
        \Rightarrow \azul{(50{,}04;\ 53{,}96)}\\[0.25em]
      99\%:\;& 52\pm 2{,}576(1{,}00)=52\pm 2{,}576
        \Rightarrow \azul{(49{,}42;\ 54{,}58)}
    \end{align*}
  \end{block}

  \vspace{-0.02cm}
  \begin{exampleblock}{Puntos 3 y 4}
    \scriptsize
    El del 99\% es más ancho: \azul{más confianza exige más margen}. Para
    estar más seguros de acertar hay que apostar a un rango mayor.

    \vspace{0.08cm}
    Interpretación: con el método usado, en el 95\% de las muestras el
    intervalo construido contendría el verdadero valor promedio de las
    facturas.
  \end{exampleblock}
\end{frame}

\begin{frame}{Ejercicio 2 — en clase}
  \begin{alertblock}{En parejas (9 minutos)}
    Se mide el tiempo de atención en \destacar{$n=20$} servicios. Resultado:
    $\bar{x}=15{,}6$ minutos y \destacar{$s=2{,}4$} minutos. Se supone
    población aproximadamente normal.

    \vspace{0.25cm}
    \begin{enumerate}
      \item ¿Se usa $z$ o $t$? ¿Con cuántos grados de libertad?
      \item Construya el intervalo del 95\% para $\mu$.
      \item ¿Qué pasaría con el ancho si $n$ fuera 80 en vez de 20?
    \end{enumerate}
  \end{alertblock}
\end{frame}

\begin{frame}{Ejercicio 2 — solución}
  \small
  \begin{block}{Se usa $t$ porque el dato es $s$, no $\sigma$}
    \[ gl=n-1=19 \qquad t_{0{,}025;\,19}=2{,}093
       \qquad \frac{s}{\sqrt{n}}=\frac{2{,}4}{\sqrt{20}}=0{,}5367 \]
    \[ 15{,}6 \pm 2{,}093(0{,}5367)=15{,}6\pm 1{,}123
       \Rightarrow \azul{(14{,}48;\ 16{,}72)} \]
  \end{block}

  \vspace{0.1cm}
  \begin{exampleblock}{Punto 3}
    \footnotesize
    Con $n=80$ el error estándar bajaría a $2{,}4/\sqrt{80}=0{,}268$ y además
    $t$ se acercaría a $1{,}99$. El intervalo se reduciría casi a la mitad:
    \azul{cuadruplicar $n$ reduce el margen a la mitad}.
  \end{exampleblock}
\end{frame}

\begin{frame}{Ejercicio 3 — en clase}
  \begin{alertblock}{Individual (8 minutos)}
    En una encuesta a \destacar{$n=400$} clientes, \destacar{248} se declaran
    satisfechos.

    \vspace{0.25cm}
    \begin{enumerate}
      \item Calcule la estimación puntual $\hat{p}$.
      \item Verifique las condiciones para usar la aproximación normal.
      \item Construya el intervalo del 95\% para $p$.
      \item ¿Puede afirmarse que más del 55\% de los clientes está
            satisfecho?
    \end{enumerate}
  \end{alertblock}
\end{frame}

\begin{frame}{Ejercicio 3 — solución}
  \begin{columns}[T,onlytextwidth]
    \begin{column}{0.52\textwidth}
      \begin{block}{Cálculos}
        \footnotesize
        \[ \hat p=\frac{248}{400}=\azul{0{,}62} \]
        \[ n\hat p=248\ge 5,\quad n(1-\hat p)=152\ge 5 \;\checkmark \]
        \[ \sqrt{\tfrac{0{,}62(0{,}38)}{400}}=0{,}0243 \]
        \[ 0{,}62\pm 1{,}96(0{,}0243)=\azul{(0{,}5724;\ 0{,}6676)} \]
      \end{block}
    \end{column}
    \begin{column}{0.44\textwidth}
      \begin{exampleblock}{Punto 4}
        \footnotesize
        Sí. Todo el intervalo está \azul{por encima} de 0{,}55, de modo que al
        95\% de confianza se puede afirmar que la satisfacción supera el
        55\%.
      \end{exampleblock}
      \begin{alertblock}{Idea clave}
        \footnotesize
        Comparar un intervalo con un valor de referencia ya es, en esencia,
        \destacar{una prueba de hipótesis}. Es lo que formalizaremos en la
        unidad 3.
      \end{alertblock}
    \end{column}
  \end{columns}
\end{frame}

\begin{frame}[fragile]{¿Qué significa ``95\% de confianza''? Véalo en R}
  \scriptsize
\begin{lstlisting}[style=Rmini]
set.seed(2026)
mu <- 100; sigma <- 15; n <- 25
cubre <- replicate(1000, {
  x  <- rnorm(n, mu, sigma)                 # una muestra nueva
  ic <- mean(x) + c(-1, 1) * qt(0.975, n - 1) * sd(x) / sqrt(n)
  ic[1] <= mu & mu <= ic[2]                 # el intervalo contuvo a mu?
})
mean(cubre)     # proporcion de intervalos que atraparon a mu
\end{lstlisting}
\begin{lstlisting}[style=Rout]
[1] 0.949
\end{lstlisting}

  \vspace{-0.05cm}
  \begin{alertblock}{La lectura correcta del 95\%}
    \scriptsize
    De 1\,000 muestras, \destacar{949 intervalos contuvieron el verdadero
    $\mu$}: cerca del 95\% prometido. La confianza es una propiedad
    \azul{del procedimiento repetido}, no de un intervalo en particular.
    Un intervalo ya calculado \destacar{contiene o no contiene} a $\mu$: no hay
    probabilidad en juego.
  \end{alertblock}

  \vspace{0.05cm}
  {\tiny\color{uniGris} Este es el experimento que aclara de una vez la
  confusión más frecuente de toda la unidad.}
\end{frame}

\begin{frame}{Cierre de la sesión}
  \begin{block}{Lo que debe quedar claro hoy}
    \begin{itemize}
      \item Un estimador puntual da un número; un intervalo da un rango con
            nivel de confianza.
      \item Estructura: estimador $\pm$ valor crítico $\times$ error estándar.
      \item Más confianza $\Rightarrow$ intervalo más ancho. Más $n$
            $\Rightarrow$ más angosto.
      \item La confianza es una propiedad del método, no del intervalo
            particular.
    \end{itemize}
  \end{block}

  \vspace{0.2cm}
  \begin{alertblock}{Trabajo independiente para la próxima sesión}
    \footnotesize
    Resolver los \destacar{20 ejercicios de la lámina siguiente}. En la sesión
    11 practicaremos el cálculo de intervalos, sus aplicaciones y sus
    limitaciones, y habrá repaso para el \destacar{Examen 2}.
  \end{alertblock}
\end{frame}

\begin{frame}{Ejercicios propuestos}
  \vspace{-0.1cm}
  \begin{beamercolorbox}[wd=\textwidth,sep=0.10cm]{block body}
    \centering\scriptsize
    $z$: 1{,}645 (90\%), 1{,}96 (95\%), 2{,}576 (99\%). \quad
    $t_{0{,}025}$: 2{,}093 ($gl$=19), 2{,}145 ($gl$=14), 2{,}045 ($gl$=29).
  \end{beamercolorbox}
  \listatiny
  \vspace{0.15cm}
  \begin{columns}[T,onlytextwidth]
    \begin{column}{0.47\textwidth}
      \begin{enumerate}\setlength{\itemsep}{0.16cm}
        \item Estimador puntual de $\mu$
        \item Estimador puntual de $p$
        \item ¿Por qué $s^{2}$ divide entre $n-1$?
        \item $\bar x=52$, $\sigma=8$, $n=64$: margen al 95\%
        \item El IC del punto 4
        \item El mismo, al 99\%
        \item El mismo, al 90\%
        \item Con $n=256$ en vez de 64: margen al 95\%
        \item ¿Qué le pasa al ancho si sube la confianza?
        \item ¿Qué le pasa al ancho si sube $n$?
      \end{enumerate}
    \end{column}
    \begin{column}{0.50\textwidth}
      \begin{enumerate}\setlength{\itemsep}{0.16cm}
        \setcounter{enumi}{10}
        \item $\bar x=15{,}6$, $s=2{,}4$, $n=20$: ¿$z$ o $t$?
        \item Grados de libertad del punto 11
        \item El IC del 95\% del punto 11
        \item $\hat p=248/400$: valor de $\hat p$
        \item ¿Se cumplen las condiciones de normalidad?
        \item Error estándar de $\hat p$
        \item IC del 95\% para $p$
        \item ¿Se puede afirmar que $p>0{,}55$?
        \item Interpretación correcta de ``95\% de confianza''
        \item ¿La confianza es del intervalo o del método?
      \end{enumerate}
    \end{column}
  \end{columns}
\end{frame}

\begin{frame}{Respuestas}{Autoevaluación de los ejercicios propuestos}
  \listatiny
  \vspace{0.15cm}
  \begin{columns}[T,onlytextwidth]
    \begin{column}{0.47\textwidth}
      \begin{enumerate}\setlength{\itemsep}{0.16cm}
        \item $\bar x$
        \item $\hat p$
        \item Para que sea insesgado
        \item $1{,}96$
        \item $(50{,}04;\ 53{,}96)$
        \item $(49{,}42;\ 54{,}58)$
        \item $(50{,}36;\ 53{,}65)$
        \item $0{,}98$ \;{\color{uniTitulo}la mitad}
        \item Se ensancha
        \item Se angosta
      \end{enumerate}
    \end{column}
    \begin{column}{0.50\textwidth}
      \begin{enumerate}\setlength{\itemsep}{0.16cm}
        \setcounter{enumi}{10}
        \item $t$, porque el dato es $s$
        \item $gl=19$
        \item $(14{,}48;\ 16{,}72)$
        \item $0{,}62$
        \item Sí: $248$ y $152$
        \item $0{,}0243$
        \item $(0{,}5724;\ 0{,}6676)$
        \item Sí: todo el IC supera $0{,}55$
        \item El 95\% de los intervalos así construidos contiene a $\mu$
        \item Del método
      \end{enumerate}
    \end{column}
  \end{columns}
\end{frame}

\begin{frame}{Referencias}
  \footnotesize
  \begin{itemize}
    \item[{[1]}] Díaz Rodríguez, M. (2022). \textit{Estadística inferencial
      aplicada}. Universidad del Norte.
      \texttt{https://elibro.net/es/lc/uniminuto/titulos/221614}
    \item[{[2]}] Llinás Solano, H. (2014). \textit{Introducción a la
      estadística con aplicaciones en ciencias sociales}. Universidad del
      Norte. \texttt{https://elibro.net/es/lc/uniminuto/titulos/70062}
    \item[{[3]}] Contento Rubio, M. R. (2019). \textit{Estadística con
      aplicaciones en R}. Editorial Utadeo.
      \texttt{https://elibro.net/es/lc/uniminuto/titulos/220926}
  \end{itemize}

  \vspace{0.35cm}
  \begin{center}
    {\Large\color{uniAzul}\textbf{¿Preguntas?}}
  \end{center}
\end{frame}

\end{document}
